\chapter{Time Series}\label{ch:iv1:time-series}

\begin{ivtrains}
The working tools of return and volatility modelling: stationarity, AR and MA, GARCH, range and high-frequency volatility estimators, unit roots and cointegration, and how to compare forecasts. Questions ask for a number and for the reason that the number misleads.
\end{ivtrains}

\begin{question}[$\star$ \iqroles{QR, QT, DS, Risk} \iqfirm{any} \iqlevel{junior}]\label{iq:iv1:time-series:1}
Which of these are weakly stationary? (a) White noise. (b) A random walk. (c) An AR(1) with coefficient $0.9$, started in its stationary distribution. (d) A linear trend plus noise. (e) The first difference of a random walk. (f) $x_t = A\cos(t) + B\sin(t)$, where $A$ and $B$ are independent standard normals, drawn once.
\end{question}

\begin{question}[$\star$ \iqroles{QR, QT, DS, Risk} \iqfirm{any} \iqlevel{junior}]\label{iq:iv1:time-series:2}
An AR(1) has $x_t = 0.7x_{t-1} + \varepsilon_t$, with noise variance 1. What are the variance and the autocorrelation at lag 3? What is the half-life of a shock? How does the $R^2$ of the $h$-step forecast depend on $h$, and what is it at $h = 1$ and $h = 2$?
\end{question}

\begin{question}[$\star$ \iqroles{QR, QT, DS, Risk} \iqfirm{any} \iqlevel{junior}]\label{iq:iv1:time-series:3}
An MA(1) has $x_t = \varepsilon_t + 0.5\varepsilon_{t-1}$, with noise variance 1. Give the variance and the autocorrelations at lags 1 and 2. What is the largest possible lag-1 autocorrelation of an MA(1)? Which other value of the coefficient gives the same autocorrelations?
\end{question}

\begin{question}[$\star$ \iqroles{QR, QT, DS, Risk} \iqfirm{any} \iqlevel{junior}]\label{iq:iv1:time-series:4}
An exponentially weighted variance uses a decay of $0.94$ on daily squared returns. What weight does the latest day get? What is the half-life of the weights, and the effective number of days? A large move of 5 times the current daily volatility arrives. By what factor does the volatility estimate jump?
\end{question}

\begin{question}[$\star\star$ \iqroles{QR, QT, DS} \iqfirm{any} \iqlevel{junior, experienced}]\label{iq:iv1:time-series:5}
Consider $x_t = 1.2x_{t-1} - 0.5x_{t-2} + \varepsilon_t$. Is it stationary? Find the roots of the characteristic polynomial, the modulus and the period of any cycle, and the autocorrelations at lags 1 and 2. What would a shock's effect look like over time?
\end{question}

\begin{question}[$\star\star$ \iqroles{QR, QT, DS} \iqfirm{any} \iqlevel{junior}]\label{iq:iv1:time-series:6}
A series of 500 points has sample autocorrelations $0.60$, $0.35$, $0.22$, $0.12$ and sample partial autocorrelations $0.60$, $-0.02$, $0.01$, $0.03$. What model do you propose, and what bands do you use to judge significance? What pattern would an MA(1) show, and what would an ARMA(1,1) show?
\end{question}

\begin{question}[$\star\star$ \iqroles{QR, QT, DS} \iqfirm{hedge fund, bank} \iqlevel{junior, experienced}]\label{iq:iv1:time-series:7}
Returns follow an AR(1) with autocorrelation $0.1$. What is the variance ratio $\operatorname{Var}(r_t + \dots + r_{t-q+1})/(q\operatorname{Var}r_t)$ for $q = 2$ and $q = 5$? What is the limit for large $q$? How much does the usual square-root-of-time annualisation understate volatility?
\end{question}

\begin{question}[$\star\star$ \iqroles{QR, QT, DS, Risk} \iqfirm{hedge fund, bank} \iqlevel{junior, experienced}]\label{iq:iv1:time-series:8}
A GARCH(1,1) has $\sigma_t^2 = \omega + \alpha r_{t-1}^2 + \beta\sigma_{t-1}^2$, with $\alpha = 0.1$, $\beta = 0.85$ and $\omega = 2 \times 10^{-6}$. Find the long-run variance and the annual volatility, with 252 days. How long does a volatility shock take to halve? If today's variance is four times the long-run value, what is the forecast 10 days ahead?
\end{question}

\begin{question}[$\star\star$ \iqroles{QR, QT, DS, Risk} \iqfirm{hedge fund, market maker} \iqlevel{junior, experienced}]\label{iq:iv1:time-series:9}
A day's high and low have ratio $H/L = 1.02$. What does the Parkinson estimator $(\ln(H/L))^2/(4\ln 2)$ give as a daily volatility? How much more efficient is it than one squared close-to-close return? Give the two main ways in which it goes wrong on real data.
\end{question}

\begin{question}[$\star\star$ \iqroles{QR, QT} \iqfirm{market maker, hedge fund} \iqlevel{junior, experienced}]\label{iq:iv1:time-series:10}
The observed log price is the efficient price plus independent noise of standard deviation $0.05\%$, and the efficient price has a daily volatility of 1\%. What does the realised variance from $n$ equal intraday returns have as its mean? What is the relative bias with 78 returns a day, and with 390? How does the noise show up in the autocovariance of returns?
\end{question}

\begin{question}[$\star\star$ \iqroles{QR, QT, DS} \iqfirm{hedge fund, bank} \iqlevel{junior, experienced}]\label{iq:iv1:time-series:11}
You estimate $x_t = \rho x_{t-1} + \varepsilon_t$ when in truth $\rho = 1$. Why does the usual $t$-statistic on $\hat{\rho} - 1$ not use the normal table? What are the 5\% critical values with no constant and with a constant? With a constant and 250 observations, about what value of $\hat{\rho}$ do you expect? Show these by simulation.
\end{question}

\begin{question}[$\star\star$ \iqroles{QR, QT, DS} \iqfirm{hedge fund} \iqlevel{junior, experienced}]\label{iq:iv1:time-series:12}
Two forecasts are compared over 400 periods. The difference in their squared errors has mean $0.02$ and standard deviation $0.2$. Find the Diebold--Mariano statistic if these differences are serially independent, and if they are AR(1) with coefficient $0.2$. What changes if one model was chosen as the best of many?
\end{question}

\begin{question}[$\star\star$ \iqroles{QR, QT} \iqfirm{hedge fund} \iqlevel{junior, experienced}]\label{iq:iv1:time-series:13}
$x_t$ is a random walk, and $y_t = 2x_t + u_t$, where $u_t$ is an AR(1) with coefficient $0.9$. What are the orders of integration of $x$, $y$ and $u$? What are the half-life of the spread, and the rate at which the regression slope converges? Describe the two steps of Engle--Granger, and say why the second step cannot use the usual table. How does this differ from regressing two independent random walks?
\end{question}

\begin{question}[$\star\star\star$ \iqroles{QR, QT, DS} \iqfirm{hedge fund, bank} \iqlevel{experienced}]\label{iq:iv1:time-series:14}
For the GARCH(1,1) with normal innovations, $\alpha = 0.1$ and $\beta = 0.85$, the kurtosis is $3(1 - (\alpha + \beta)^2)/(1 - (\alpha + \beta)^2 - 2\alpha^2)$. Derive this by taking the fourth moment. Evaluate it. What is the condition for it to be finite? The returns themselves are uncorrelated. What are the lag-1 autocorrelation of the squared returns, given by $\alpha(1 - \alpha\beta - \beta^2)/(1 - 2\alpha\beta - \beta^2)$, and its decay rate?
\end{question}

\begin{question}[$\star\star\star$ \iqroles{QR, QT, DS} \iqfirm{hedge fund, bank} \iqlevel{experienced}]\label{iq:iv1:time-series:15}
A perfect forecast of the conditional variance $\sigma_t^2$ is scored against the squared return $r_t^2$. Show that the $R^2$ of this regression is at most $(\kappa/3 - 1)/(\kappa - 1)$, where $\kappa$ is the return's kurtosis and the innovations are normal. Evaluate this for the GARCH(1,1) with kurtosis $3.77$. What does it say about judging volatility models? How do you fix the comparison?
\end{question}

\begin{question}[$\star\star\star$ \iqroles{QR, QT} \iqfirm{hedge fund} \iqlevel{experienced}]\label{iq:iv1:time-series:16}
A return is $r_t = \mu_t + \varepsilon_t$, where the expected return $\mu_t$ is a slow AR(1) with coefficient $\phi = 0.9$ and where the noise is white and independent of it. The signal is $R^2 = 1\%$ of total variance. What is the lag-1 autocorrelation of returns, and at lag $k$? What ARMA model does $r$ follow? Why can return autocorrelations look like nothing while the predictability is real?
\end{question}

\begin{question}[$\star\star\star$ \iqroles{QR, QT, DS, Risk} \iqfirm{hedge fund, bank} \iqlevel{experienced}]\label{iq:iv1:time-series:17}
Returns are normal given a state, with annualised volatility of 10\% in a calm state and 30\% in a stressed state. The state is a Markov chain that stays calm with probability $0.95$ and stressed with probability $0.9$. Find the stationary shares, the unconditional volatility and the kurtosis. How does the autocorrelation of squared returns decay?
\end{question}

\begin{question}[$\star\star\star$ \iqroles{QR, QT, DS} \iqfirm{hedge fund} \iqlevel{junior, experienced}]\label{iq:iv1:time-series:18}
You fit an AR(1) by least squares with a constant to 100 points of a mean-reverting spread, where the true coefficient is $0.95$. Use the bias formula $\mathbb{E}\hat{\phi} - \phi \approx -(1 + 3\phi)/T$ to find the expected estimate. What does that do to the half-life you report? Why does the bias arise, and how can you correct it?
\end{question}

\begin{question}[$\star\star\star\star$ \iqroles{QR, QT} \iqfirm{hedge fund} \iqlevel{experienced}]\label{iq:iv1:time-series:19}
Under a unit root with no constant, derive the limit laws of $T(\hat{\rho} - 1)$ and of the $t$-statistic, as functionals of a Brownian motion on $[0, 1]$. Find the distribution of the numerator, and the probability that it is negative. Check the 5\% critical value of the $t$-statistic and the mean of $T(\hat{\rho} - 1)$ by simulation.
\end{question}

\begin{question}[$\star\star\star\star$ \iqroles{QR, QT} \iqfirm{market maker, hedge fund} \iqlevel{experienced}]\label{iq:iv1:time-series:20}
The efficient price has daily variance $\sigma^2$, and the noise has variance $\eta^2$, as in a standard model of market microstructure noise. Write the mean squared error of realised variance as a function of $n$, keeping the main terms, and find the best $n$, for $\sigma = 1\%$ and $\eta = 0.05\%$. Then show how averaging over subsamples and subtracting a bias term gives an estimator that uses all the data and has no first-order bias, and test it by simulation.
\end{question}
